\subsection{\texorpdfstring{THE PLANE \(E_{s,s'}\) AND THE GROUP GENERATED BY \(\sigma_{s}\) AND \(\sigma_{s'}\)}{THE PLANE E\_s,s' AND THE GROUP GENERATED BY sigma\_s AND sigma\_s'}}

In this number, we denote by \(s\) and \(s'\) two elements of \(S\) with \(s \neq s'\) . Put \(m = m(s, s')\) and denote by \(\mathrm{E}_{s,s'}\) the plane \(\mathbf{Re}_s \oplus \mathbf{Re}_{s'}\) .

PROPOSITION 1. The restriction of \(B_{M}\) to \(E_{s,s'}\) is positive, and it is nondegenerate if and only if \(m\) is finite.

Let \(z = xe_{s} + ye_{s'}\) , with \(x, y \in \mathbf{R}\) , be an element of \(E_{s,s'}\) . We have

\[
\begin{array}{r} B _ {M} (z, z) = x ^ {2} - 2 x y \cos \frac {\pi}{m} + y ^ {2} \\ = (x - y. \cos \frac {\pi}{m}) ^ {2} + y ^ {2} \sin^ {2} \frac {\pi}{m}, \end{array}
\]

which shows that \(\mathrm{B}_{\mathsf{M}}\) is positive on \(\mathrm{E}_{s,s'}\) , and that it is non-degenerate there if and only if \(\sin \frac{\pi}{m} \neq 0\) . The proposition follows.

The reflections \(\sigma_{s}\) and \(\sigma_{s'}\) leave \(E_{s,s'}\) stable. We are going to determine the order of the restriction of \(\sigma_{s}\sigma_{s'}\) to \(E_{s,s'}\) . We distinguish two cases:

\begin{enumerate}
\def\labelenumi{\alph{enumi})}
\tightlist
\item
  \(m = +\infty\)
\end{enumerate}

Let \(u = e_s + e_{s'}\) . We have \(\mathrm{B}_{\mathrm{M}}(u, e_s) = \mathrm{B}_{\mathrm{M}}(u, e_{s'}) = 0\) , so \(u\) is invariant under \(\sigma_s\) and \(\sigma_{s'}\) . Moreover,

\[
\sigma_ {s} \left(\sigma_ {s ^ {\prime}} \left(e _ {s}\right)\right) = \sigma_ {s} \left(e _ {s} + 2 e _ {s ^ {\prime}}\right) = 3 e _ {s} + 2 e _ {s ^ {\prime}} = 2 u + e _ {s},
\]

hence

\[
\left(\sigma_ {s} \sigma_ {s ^ {\prime}}\right) ^ {n} \left(e _ {s}\right) = 2 n u + e _ {s} \text {   for   all   } n \in \mathbf {Z}.
\]

It follows that the restriction of \(\sigma_s\sigma_{s'}\) to \(\mathrm{E}_{s,s'}\) is of infinite order.

\begin{enumerate}
\def\labelenumi{\alph{enumi})}
\setcounter{enumi}{1}
\tightlist
\item
  \(m\) is finite.
\end{enumerate}

The form \(B_{M}\) provides \(E_{s,s'}\) with the structure of a euclidean plane. Since the scalar product of \(e_{s}\) and \(e_{s'}\) is equal to \(-\cos\frac{\pi}{m}=\cos\left(\pi-\frac{\pi}{m}\right)\) , we can orient \(E_{s,s'}\) so that the angle between the half-lines \(R_{+}e_{s}\) and \(R_{+}e_{s'}\) is equal to \(\pi-\frac{\pi}{m}\) . If D and \(D'\) denote the lines orthogonal to \(e_{s}\) and \(e_{s'}\) ,

\[
(\widehat {D ^ {\prime} , D}) = \pi - (\widehat {D , D ^ {\prime}}) = \frac {\pi}{m}.
\]

Now, the restrictions \(\overline{\sigma}_s\) and \(\overline{\sigma}_{s'}\) of \(\sigma_s\) and \(\sigma_{s'}\) to \(E_{s,s'}\) are the orthogonal symmetries with respect to D and \(D'\) . By the Cor. to Prop. 6 of §2, no. 5, it follows that \(\overline{\sigma}_s\overline{\sigma}_{s'}\) is the rotation with angle \(\frac{2\pi}{m}\) . In particular, its order is \(m\) .

We now return to E:

PROPOSITION 2. The subgroup of \(\mathbf{GL}(\mathrm{E})\) generated by \(\sigma_s\) and \(\sigma_{s'}\) is a dihedral group of order \(2m(s,s')\) .

Since \(\sigma_s\) and \(\sigma_{s'}\) are of order 2, and are distinct, it is enough to show that their product \(\sigma_s\sigma_{s'}\) is of order \(m(s,s')\) . When \(m(s,s')\) is infinite, this follows from a) above. When \(m(s, s')\) is finite, it follows from Prop. 1 that E is the direct sum of \(E_{s,s'}\) and its orthogonal complement \(V_{s,s'}\) ; since \(\sigma_s\) and \(\sigma_{s'}\) are the identity on \(V_{s,s'}\) , and since the restriction of \(\sigma_s\sigma_{s'}\) to \(E_{s,s'}\) is of order \(m(s, s')\) by b), the order of \(\sigma_s\sigma_{s'}\) is indeed equal to \(m(s, s')\) .

